\documentclass{article}
\usepackage[utf8]{inputenc}
\usepackage{a4wide}

\begin{document}

\section*{สวนฉันจะไปประกวดต่างดาว}

เรนิตาเป็นผู้ดูแลสวน\texttt{n}สวนและต้องการที่จะปลูกดอกไม้เพื่อตกแต่งสวน



กำหนดให้0-indexedinteger arrayขนาด\texttt{n}ชื่อ\texttt{flowers}แทนดอกไม้

โดยที่\texttt{flowers{[}i{]}}แทนจำนวนของดอกไม้ที่ปลูกแล้วในสวนที่\texttt{i}ดอกไม้ที่ที่ปลูกไปแล้วไม่สามารถถอนออกได้



คุณได้มอบดอกไม้ใหม่

แทนด้วยตัวแปร\texttt{newFlowers}ซึ่งหมายถึงจำนวนมากที่สุดของดอกไม้ที่เรนิตาสามารถปลูกเพิ่มได้และคุณได้กำหนดค่าเพิ่มอีก\texttt{3}ค่า

ได้แก่target, fullและpartial



สวนหนึ่งๆ จะสมบูรณ์เมื่อมีดอกไม้ถูกปลูกอย่างน้อยเท่าtarget



ค่าความสวยรวมของสวนคิดได้จากผลรวมของค่าดังต่อไปนี้



\begin{enumerate}

\tightlist

\item

  จำนวนสวนที่สมบูรณ์คูณด้วยfullหมายถึงจำนวน

\item

  จำนวนน้อยที่สุดของดอกไม้ที่ปลูกอยู่ในสวนที่ไม่สมบูรณ์ในกรณีที่ไม่มีสวนที่ไม่สมบูรณ์ค่านี้จะมีค่าเป็น0

\end{enumerate}



จงเขียนโปรแกรมที่\textbf{คืนค่าความสวยรวมสูงที่สุด}ที่เรนิตาจะสามารถทำได้หลังจากปลูกดอกไม้ให้ได้มากที่สุด



\subsection*{Input}
จำนวนสวนของเรนิตาทั้งหมด\texttt{n}สวน



จำนวนดอกไม้ในแต่ละสวนของเรนิตา\texttt{flowersi}



ดอกไม้ใหม่ที่เรนิตาสามารถปลูกเพิ่มได้\texttt{newFlowers}และค่าของ\texttt{target,\ full,\ partial}



\subsection*{Output}
ความสวยรวมสูงที่สุดของสวนหลังจากปลูกดอกไม้ให้ได้มากที่สุด\\



\end{document}
